% Part: sets-functions-relations
% Chapter: sets
% Section: unions-and-intersections

\documentclass[../../../include/open-logic-section]{subfiles}

\begin{document}

\olfileid{sfr}{set}{uni}
\olsection{并集与交集}

\begin{explain}
在 \olref[sfr][set][bas]{sec} 中，我们引入了通过抽象给出的集合定义，即形如 $\Setabs{x}{\phi(x)}$ 的定义。这里，我们使用了某个性质~$\phi$，而这个性质可以提及我们已经定义过的集合。举例来说，如果 $A$ 和~$B$ 是集合，那么集合 $\Setabs{x}{x \in A \lor x \in B}$ 由所有属于 $A$ 或~$B$ 的对象组成，也就是说，它是合并了 $A$ 和~$B$ 的元素的集合。我们可以如 \olref{fig:union} 那样将其可视化，其中高亮区域表示两个集合 $A$ 和~$B$ 合在一起的元素。

\begin{figure}
  \olcachedasset{union}
  \caption{两个集合的并集 $A \cup B$ 是 $A$ 的元素与~$B$ 的元素共同组成的集合。}
  \ollabel{fig:union}
\end{figure}

这种对集合的操作——将它们合并——非常有用且常见，因此我们给它一个正式的名称和符号。
\end{explain}

\begin{defn}[并集]
两个集合 $A$ 和 $B$ 的\emph{并集}，记作 $A \cup B$，是由所有属于 $A$、属于 $B$ 或同时属于二者的元素构成的集合。
\[
A \cup B = \Setabs{x}{x \in A \lor x \in B}
\]
\end{defn}

\begin{ex}
由于元素的重数无关紧要，有共同元素的两个集合的并集中，该元素只出现一次，例如，$\{ a, b, c\} \cup \{ a, 0, 1\} = \{a, b, c, 0, 1\}$。

一个集合与其某个子集的并集就是那个较大的集合：$\{a,
b, c \} \cup \{a \} = \{a, b, c\}$。

一个集合与空集的并集等于该集合本身：$\{a,
b, c \} \cup \emptyset = \{a, b, c \}$。
\end{ex}

\begin{prob}
证明：若 $A \subseteq B$，则 $A \cup B = B$。
\end{prob}

\begin{explain}
我们也可以考虑一种与并集“对偶”的操作。这种操作生成由所有既属于~$A$ 又属于~$B$ 的元素构成的集合。这个操作称为\emph{交集}，可以像 \olref{fig:intersection} 中那样描绘。
\begin{figure}
  \olcachedasset{intersection}
  \caption{两个集合的交集 $A \cap B$ 是它们共同拥有的元素组成的集合。}
  \ollabel{fig:intersection}
\end{figure}
\end{explain}

\begin{defn}[交集]
两个集合 $A$ 和 $B$ 的\emph{交集}，记作 $A \cap B$，是由所有同时属于 $A$ 和~$B$ 的元素构成的集合。
\[
A \cap B = \Setabs{x}{x \in A \land x \in B}
\]
若两个集合的交集为空，则称它们是\emph{不相交}的。这意味着它们没有共同的元素。
\end{defn}

\begin{ex}
如果两个集合没有共同的元素，它们的交集就是空集：$\{ a, b, c\} \cap \{ 0, 1\} = \emptyset$。

如果两个集合确实有共同的元素，它们的交集就是所有这些元素的集合：$\{a, b, c \} \cap \{a, b, d \} = \{a, b\}$。

一个集合与其某个子集的交集就是那个较小的集合：$\{a, b, c\} \cap \{a, b\} = \{a, b\}$。

空集与任何集合的交集都是空集：$\{a, b, c \}
\cap \emptyset = \emptyset$
\end{ex}

\begin{prob}
严格证明：若 $A \subseteq B$，则 $A \cap B = A$。
\end{prob}

\begin{explain}
我们也可以对多于两个的集合作并或交。一种一般且优雅的处理方式如下：假设你将所有想要取并集（或交集）的集合收集到一个单独的集合中。那么，我们可以将所有原始集合的并集定义为属于该集合至少一个元素的对象所组成的集合，而将交集定义为属于该集合每一个元素的对象所组成的集合。
\end{explain}

\begin{defn}
如果 $A$ 是一个由集合构成的集合，那么 $\bigcup A$ 是 $A$ 中元素的元素所组成的集合：
\begin{align*}
\bigcup A & = \Setabs{x}{x \text{ 属于 } A},
\text{ 的某个元素，即，}\\
& = \Setabs{x}{\text{存在 } B \in A
  \text{ 使得 } x \in B}
\end{align*}
\end{defn}

\begin{defn}
如果 $A$ 是一个由集合构成的集合，那么 $\bigcap A$ 是 $A$ 中所有元素共有的对象所组成的集合：
\begin{align*}
\bigcap A & = \Setabs{x}{x \text{ 属于 } A},
\text{ 的每一个元素，即，}\\
 & = \Setabs{x}{\text{ 对所有 } B \in A, x \in B}
\end{align*}
\end{defn}

\begin{ex}
假设 $A = \{ \{ a, b \}, \{ a, d, e \}, \{ a, d \} \}$。
那么 $\bigcup A = \{ a, b, d, e \}$，且 $\bigcap A = \{ a \}$。
\end{ex}
\begin{prob}
	证明：如果 $A$ 是一个集合且 $A \in B$，那么 $A \subseteq \bigcup B$。
\end{prob}

我们也可以对集合序列 $A_1$, $A_2$, \dots 作同样的处理：
\begin{align*}
\bigcup_i A_i & = \Setabs{x}{x \text{ 属于某个 } A_i}\\
\bigcap_i A_i & = \Setabs{x}{x \text{ 属于每一个 } A_i}.
\end{align*}

当我们有集合的\emph{指标}时，即存在某个集合 $I$，使得我们对于每个 $i \in I$ 考虑 $A_i$，我们也可以使用这些缩写：
\begin{align*}
	\bigcup_{i \in I} A_i & = \bigcup \Setabs{A_i }{i \in I}\\
	\bigcap_{i \in I} A_i & = \bigcap\Setabs{A_i}{i \in I}
\end{align*}

最后，我们可能要考虑 $A$ 中那些不在 $B$ 里的所有元素组成的集合。我们可以如 \olref{difference} 那样描绘它。

\begin{figure}
  \olcachedasset{difference}
  \caption{两个集合的差 $A \setminus B$ 是 $A$ 中那些不属于 $B$ 的元素组成的集合。}
  \ollabel{difference}
\end{figure}

\begin{defn}[差集]
$A$ 与 $B$ 的\emph{集合差}~$A \setminus B$ 是 $A$ 中所有不属于 $B$ 的元素组成的集合，即
\[
A\setminus B = \Setabs{x}{x\in A \text{ 且 } x \notin B}.
\]
\end{defn}

\begin{prob}
证明：如果 $A \subsetneq B$，那么 $B \setminus A \neq \emptyset$。
\end{prob}

\end{document}
